\subsection{What Is Meant by Data?}
\label{Section 2.1.3}

We began the rational-number implementation in \link{Section 2.1.1} by
implementing the rational-number operations \code{add\_rat}, \code{sub\_rat}, and
so on in terms of three unspecified procedures: \code{make\_rat}, \code{numer},
and \code{denom}.  At that point, we could think of the operations as being
defined in terms of data objects---numerators, denominators, and rational
numbers---whose behavior was specified by the latter three procedures.

But exactly what is meant by \newterm{data}?  It is not enough to say
``whatever is implemented by the given selectors and constructors.''  Clearly,
not every arbitrary set of three procedures can serve as an appropriate basis
for the rational-number implementation.  We need to guarantee that, if we
construct a rational number \code{x} from a pair of integers \code{n} and
\code{d}, then extracting the \code{numer} and the \code{denom} of \code{x} and
dividing them should yield the same result as dividing \code{n} by \code{d}.
In other words, \code{make\_rat}, \code{numer}, and \code{denom} must satisfy
the condition that, for any integer \code{n} and any non-zero integer \code{d},
if \code{x} is \code{make\_rat(n, d)}, then

\[
\frac{\texttt{numer(x)}}{\texttt{denom(x)}}
=
\frac{\texttt{n}}{\texttt{d}}\,.
\]

In fact, this is the only condition \code{make\_rat}, \code{numer}, and
\code{denom} must fulfill in order to form a suitable basis for a
rational-number representation.  In general, we can think of data as defined by
some collection of selectors and constructors, together with specified
conditions that these procedures must fulfill in order to be a valid
representation.\footnote{Surprisingly, this idea is very difficult to formulate
rigorously. There are two approaches to giving such a formulation.  One,
pioneered by C. A. R. \link{Hoare (1972)}, is known as the method of \newterm{abstract
models}.  It formalizes the ``procedures plus conditions'' specification as
outlined in the rational-number example above.  Note that the condition on the
rational-number representation was stated in terms of facts about integers
(equality and division).  In general, abstract models define new kinds of data
objects in terms of previously defined types of data objects.  Assertions about
data objects can therefore be checked by reducing them to assertions about
previously defined data objects.  Another approach, introduced by Zilles at
\acronym{MIT}, by Goguen, Thatcher, Wagner, and Wright at \acronym{IBM}
(see \link{Thatcher et al. 1978}),
and by Guttag at Toronto (see \link{Guttag 1977}), is called
\newterm{algebraic specification}.  It regards the ``procedures'' as elements
of an abstract algebraic system whose behavior is specified by axioms that
correspond to our ``conditions,'' and uses the techniques of abstract algebra
to check assertions about data objects.  Both methods are surveyed in the paper
by \link{Liskov and Zilles (1975)}.}

This point of view can serve to define not only ``high-level'' data objects,
such as rational numbers, but lower-level objects as well.  Consider the notion
of a pair, which we used in order to define our rational numbers.  We never
actually said what a pair was, only that the language supplied procedures
\code{cons}, \code{car}, and \code{cdr} for operating on pairs.  But the only
thing we need to know about these three operations is that if we glue two
objects together using \code{cons} we can retrieve the objects using \code{car}
and \code{cdr}.  That is, the operations satisfy the condition that, for any
objects \code{x} and \code{y}, if \code{z} is \code{cons(x, y)} then
\code{car(z)} is \code{x} and \code{cdr(z)} is \code{y}.  We have not troubled
to write those three procedures until now, having had the tuple to hand: what
we wrote as \code{(x, y)}, \code{z[0]} and \code{z[1]} is what the original
writes as \code{cons}, \code{car} and \code{cdr}, and the tuple, like the
original's pair, is supplied by the language and not by us.

Any triple of procedures that satisfies the above condition can be used as the
basis for implementing pairs.  This point is illustrated strikingly by the fact
that we could implement \code{cons}, \code{car}, and \code{cdr} without using
any data structure at all---without the tuple, or anything else the language
provides for holding two things together---but only procedures.  Here are the
definitions:

\begin{codeBlock}
>>> def cons(x, y):
...     def dispatch(m):
...         if m == 0:
...             return x
...         elif m == 1:
...             return y
...         else:
...             raise ValueError(f"Argument not 0 or 1: cons {m}")
...     return dispatch

>>> def car(z):
...     return z(0)

>>> def cdr(z):
...     return z(1)
\end{codeBlock}

\noindent
This use of procedures corresponds to nothing like our intuitive notion of what
data should be.  Nevertheless, all we need to do to show that this is a valid
way to represent pairs is to verify that these procedures satisfy the condition
given above.

\begin{revealTheSurface}
Intuition says that a procedure is a piece of code and not a place to keep
things.  Python will settle the question, because a closure---the word was
introduced in \link{Section 1.1.8}---carries its captured values about with it,
and lets you look:

\begin{codeBlock}
>>> z = cons(1, 2)
>>> z
~\textit{<function cons.<locals>.dispatch at 0x104727270>}~
>>> z.__code__.co_freevars
~\textit{('x', 'y')}~
>>> [cell.cell_contents for cell in z.__closure__]
~\textit{[1, 2]}~
\end{codeBlock}

\noindent
\code{co\_freevars} names the variables \code{dispatch} did not define for
itself and so captured from around it, and \code{\_\_closure\_\_} holds a
\newterm{cell} for each, which is where the value is actually kept.  The 1 and
the 2 are in there.  They are not written down anywhere in \code{dispatch}, and
\code{cons} returned long ago, but a procedure that mentions \code{x} and
\code{y} cannot be run without them, so they had to be put somewhere and they
were.

This is the footnote of \link{Section 1.3.4} made visible: giving procedures
full first-class status, we said there, means reserving storage for a
procedure's free variables even while the procedure is not executing.  The
cells are that storage.  A procedure that closes over two values is a thing
that holds two values, whatever our intuitions about code and data may
say---and that, rather than any cleverness in the dispatching, is why the
representation above works at all.
\end{revealTheSurface}

The subtle point to notice is that the value returned by \code{cons(x, y)} is a
procedure---namely the internally defined procedure \code{dispatch}, which
takes one argument and returns either \code{x} or \code{y} depending on whether
the argument is 0 or 1.  Correspondingly, \code{car(z)} is defined to apply
\code{z} to 0.  Hence, if \code{z} is the procedure formed by
\code{cons(x, y)}, then \code{z} applied to 0 will yield \code{x}.  Thus, we
have shown that \code{car(cons(x, y))} yields \code{x}, as desired.  Similarly,
\code{cdr(cons(x, y))} applies the procedure returned by \code{cons(x, y)} to
1, which returns \code{y}.  Therefore, this procedural implementation of pairs
is a valid implementation, and if we access pairs using only \code{cons},
\code{car}, and \code{cdr} we cannot distinguish this implementation from one
that uses ``real'' data structures.

The point of exhibiting the procedural representation of pairs is not that our
language works this way (Python implements the tuple directly, for efficiency
reasons, as Lisp systems implement the pair) but that it could work this way.
The procedural representation, although obscure, is a perfectly adequate way to
represent pairs, since it fulfills the only conditions that pairs need to
fulfill.  This example also demonstrates that the ability to manipulate
procedures as objects automatically provides the ability to represent compound
data.  This may seem a curiosity now, but procedural representations of data
will play a central role in our programming repertoire.  This style of
programming is often called \newterm{message passing}, and we will be using it
as a basic tool in \link{Chapter 3} when we address the issues of modeling and
simulation.  A reader who knows Python will suspect they have seen this shape
before, and will be right; we shall say plainly what it is when we get there.

\begin{exerciseGroup}
\phantomsection\label{Exercise 2.4}
\exerciseHeading{Exercise 2.4:} Here is an alternative procedural
representation of pairs.  For this representation, verify that
\code{car(cons(x, y))} yields \code{x} for any objects \code{x} and \code{y}.

\begin{codeBlock}
>>> def cons(x, y):
...     return lambda m: m(x, y)

>>> def car(z):
...     return z(lambda p, q: p)
\end{codeBlock}

What is the corresponding definition of \code{cdr}? (Hint: To verify that this
works, make use of the substitution model of \link{Section 1.1.5}.)

\phantomsection\label{Exercise 2.5}
\exerciseHeading{Exercise 2.5:} Show that we can represent pairs of
nonnegative integers using only numbers and arithmetic operations if we
represent the pair \( a \) and \( b \) as the integer that is the product \( 2^a 3^b \).
Give the corresponding definitions of the procedures \code{cons},
\code{car}, and \code{cdr}.

\phantomsection\label{Exercise 2.6}
\exerciseHeading{Exercise 2.6:} In case representing pairs as
procedures wasn't mind-boggling enough, consider that, in a language that can
manipulate procedures, we can get by without numbers (at least insofar as
nonnegative integers are concerned) by implementing 0 and the operation of
adding 1 as

\begin{codeBlock}
>>> zero = lambda f: lambda x: x

>>> def add_1(n):
...     return lambda f: lambda x: f(n(f)(x))
\end{codeBlock}

This representation is known as \newterm{Church numerals}, after its inventor,
Alonzo Church, the logician who invented the λ-calculus.

A Church numeral is a procedure, so asking the interpreter to show one tells
you only that it is a procedure.  To see which number it stands for, apply it.
A numeral applies its argument that many times to a starting value, so give it
a procedure that adds one and a starting value of zero, and the number itself
comes back:

\begin{codeBlock}
>>> def evaluate_church(n):
...     return n(lambda x: x + 1)(0)

>>> evaluate_church(zero)
~\textit{0}~
>>> evaluate_church(add_1(zero))
~\textit{1}~
>>> evaluate_church(add_1(add_1(zero)))
~\textit{2}~
\end{codeBlock}

\noindent
Notice that \code{evaluate\_church} does no counting and no repeating of its
own.  The numeral does all of it: whatever it means to be three is contained in
the numeral, and the procedure above only has to supply something to do three
times.

Define \code{one} and \code{two} directly (not in terms of \code{zero} and
\code{add\_1}).  (Hint: Use substitution to evaluate \code{add\_1(zero)}.)  Give
a direct definition of an addition procedure (not in terms of repeated
application of \code{add\_1}).
\end{exerciseGroup}

